\documentclass[12pt]{article}
\usepackage[utf8]{inputenc}
\usepackage[T1]{fontenc}
\usepackage[spanish]{babel}
\usepackage{geometry}
\usepackage{amsmath,amssymb}
\usepackage{enumitem}

\geometry{margin=2.5cm}

\newcommand{\dudoso}[1]{\textsf{#1}}

\begin{document}

\begin{center}
{\Large Transcripción cruda}\\[2mm]
{\large 11May26}
\end{center}

\section*{Foto 1}

An\'alisis Matem\'atico I.

3 ex\'amenes parciales.
Abarcando cada uno dos cap\'itulos.

\begin{enumerate}[label=\arabic*.]
    \item Im\'agenes y preim\'agenes de conjuntos bajo la acci\'on de funciones.
    \item Espacios m\'etricos.
    \item Funciones continuas.
    \item Conceptos topol\'ogicos en espacios m\'etricos.
    \item Espacios m\'etricos completos.
    \item Espacios de funciones.
\end{enumerate}

Textos:
\begin{enumerate}[label=\arabic*.]
    \item Rudin, W., \textit{Principios de An\'alisis Matem\'atico}, 3.\textsuperscript{a} edici\'on, McGraw-Hill, \dudoso{1981}.
\end{enumerate}

Consulta:
\begin{enumerate}[label=\arabic*.]
    \item Bartle, R. G., \textit{Elementos de An\'alisis Real}, Ed. Limusa.
\end{enumerate}

\section*{Foto 2}

Sean $A$ y $B$ conjuntos. Sea $f\colon A\to B$.
Sean $C\subset A$ y $D\subset B$.

La imagen directa de $C$ bajo la acci\'on de $f$ es
\[
f(C)=\{f(a)\mid a\in C\}.
\]

La imagen inversa de $D$, bajo la acci\'on de $f$, es
\[
f^{-1}(D)=\{a\in A\mid f(a)\in D\}.
\]

\section*{Foto 3}

Repetici\'on del repaso de imagen directa e inversa.

\section*{Foto 4}

\textit{Algunas propiedades:}

\begin{enumerate}[label=\arabic*.]
    \item La imagen preserva uniones arbitrarias de conjuntos.
    \[
    f\Bigl(\bigcup_\alpha C_\alpha\Bigr)=\bigcup_\alpha f(C_\alpha).
    \]
    No preserva la intersecci\'on.
    \item Las im\'agenes inversas s\'i preservan todas las operaciones de conjuntos, es decir
    \[
    f^{-1}\Bigl(\bigcup_\alpha D_\alpha\Bigr)=\bigcup_\alpha f^{-1}(D_\alpha),
    \qquad
    f^{-1}\Bigl(\bigcap_\alpha D_\alpha\Bigr)=\bigcap_\alpha f^{-1}(D_\alpha).
    \]
\end{enumerate}

\section*{Foto 5}

\[
f^{-1}\Bigl(\bigcap_\alpha D_\alpha\Bigr)=\bigcap_\alpha f^{-1}(D_\alpha),
\]
\[
f^{-1}(\varnothing)=\varnothing,\qquad f(\varnothing)=\varnothing.
\]

\textit{Ejemplos.}
\[
f(\varnothing)=\varnothing,\qquad f^{-1}(\varnothing)=\varnothing.
\]

\textit{Def.}
Sean $A$ y $B$ conjuntos.
Decimos que son equivalentes si existe una biyecci\'on $f\colon A\to B$.
Se da una ``relaci\'on de equivalencia'' entre los conjuntos.

\section*{Foto 6}

\textit{Def.}
Sea $A$ un conjunto. Su cardinalidad es la clase de equivalencia a la que pertenece.

\textit{Def.}
Un conjunto $A$ es finito si existe $n\in\mathbb N$ tal que $A$ se puede poner en biyecci\'on con $\{1,\dots,n\}$.

\textit{Def.}
Un conjunto $A$ es numerable si existe $f\colon\mathbb N\to A$ que es biyectiva.

\textit{Ejemplos.} $\mathbb Z$.

\textit{Teorema.}
Si $A$ es numerable y $E\subset A$ es infinito, entonces $E$ es numerable.

\section*{Foto 7}

\textit{Dem. del teorema.}
Dos casos.

\begin{enumerate}[label=\arabic*.]
    \item Sup\'ongase que $E_i\cap E_j=\varnothing$ si $i\neq j$. Establezcamos una biyecci\'on entre $\bigcup_{n=1}^\infty E_n$ y $\mathbb N\times\mathbb N$.
    \item Si la familia $\{E_n\}_{n=1}^\infty$ no fuera mutuamente ajena.
\end{enumerate}

\section*{Foto 8}

Construcci\'on diagonal:
\[
(1,1),(1,2),(1,3),\dots
\]
\[
(2,1),(2,2),(2,3),\dots
\]
\[
(3,1),(3,2),(3,3),\dots
\]

Reto: dar expl\'icitamente
\[
f\colon\mathbb N\to\mathbb N\times\mathbb N
\]
siguiendo el esquema de arriba.

\section*{Foto 9}

\textit{Teorema.} Si $A$ es numerable y $E_n$ es un conjunto numerable para cada $n$, entonces
\[
S=\bigcup_{n=1}^\infty E_n
\]
es numerable.

\textit{(En, palabritas: la uni\'on numerable de numerables es numerable.)}

\textit{Lema.} $\mathbb N\times\mathbb N$ es conjunto numerable.

\textit{Dem del lema.} $\mathbb N\times\mathbb N$.

\section*{Foto 10}

En un conjunto numerable.

Reto: dar expl\'icitamente
\[
f\colon\mathbb N\to\mathbb N\times\mathbb N
\]
siguiendo el esquema de arriba.

\section*{Foto 11}

\textit{Dem. del teorema.}
Dos casos:

\begin{enumerate}[label=\arabic*.]
    \item Sup\'ongase que $E_i\cap E_j=\varnothing$ si $i\neq j$.
\end{enumerate}

Establecemos una biyecci\'on entre $\bigcup_{n=1}^\infty E_n$ y $\mathbb N\times\mathbb N$.

Sea
\[
E_n=\{x_{n1},x_{n2},x_{n3},\dots\}.
\]
Definimos
\[
T:\mathbb N\times\mathbb N\to \bigcup_{n=1}^\infty E_n,\qquad T((i,j))=x_{ij}.
\]

\section*{Foto 12}

Si $i_1\neq i_2$ entonces $T(i_1,j_1)\neq T(i_2,j_2)$ porque
\[
T(i_1,j_1)\in E_{i_1},\qquad T(i_2,j_2)\in E_{i_2},
\]
y como $E_{i_1}\cap E_{i_2}=\varnothing$, entonces
\[
T(i_1,j_1)\neq T(i_2,j_2).
\]

Si $i_1=i_2$, entonces $j_1\neq j_2$ aunque
\[
T(i_1,j_1),T(i_2,j_2)\in E_{i_1},
\]
son elementos distintos del mismo conjunto.

\section*{Foto 13}

\[
\therefore\ T(i_1,j_1)\neq T(i_2,j_2).
\]

Si la familia $\{E_n\}_{n=1}^\infty$ no fueran mutuamente ajenos, se listan los elementos
\[
x_{11}\ x_{12}\ x_{13}\ \dots
\]
\[
x_{21}\ x_{22}\ x_{23}\ \dots
\]
\[
x_{31}\ x_{32}\ x_{33}\ \dots
\]
y
\[
\bigcup_{j=1}^\infty E_j
\]
del arreglo de arriba, que es numerable.

\section*{Foto 14}

\textit{Corolario.} Si $A$ es conjunto numerable y $B$ es conjunto numerable, entonces $A\times B$ es conjunto numerable.

\textit{Dem.}
Sea
\[
A_\alpha=\{(\alpha,\beta)\mid \beta\in B\}\subset A\times B.
\]
Entonces $A_\alpha$ es numerable y
\[
A\times B=\bigcup_{\alpha\in A}A_\alpha.
\]

\textit{Corolario.} Si $n\in\mathbb N$ y $A_1,\dots,A_n$ son conjuntos numerables, entonces
\[
A_1\times\cdots\times A_n
\]
es numerable.

\textit{Lema.} $\mathbb N\times\mathbb N$ es numerable ya.

\section*{Foto 15}

\textit{Teorema.} $\mathbb Q$ es numerable.

\[
\mathbb Q=\left\{\frac{a}{b}\ \middle|\ a,b\in\mathbb Z,\ b>0\right\}.
\]

\[
\mathbb Z\times\mathbb N \overset{S}{\to}\mathbb Q,\qquad S(a,b)=\frac{a}{b}.
\]
Es suprayectiva.
\[
\mathbb Q \text{ es numerable.}
\]

\section*{Foto 16}

\textit{Ejercicio.}
Sea $A$ numerable, $B$ un conjunto, $f\colon A\to B$ suprayectiva, entonces $B$ es a lo m\'as numerable.

\textit{Ejemplo de un conjunto que es NO numerable, es decir no es finito y tampoco es numerable.}
\[
\mathcal C:=\{f:\mathbb N\to\{0,1\}\}.
\]
Entonces $\mathcal C$ es no numerable.

\section*{Foto 17}

Sea $E\subset \mathcal C$ cualquier subconjunto numerable de $\mathcal C$. P.D. que $E$ es un subconjunto propio.

\section*{Foto 18}

\textit{Ejemplo de un conjunto que es NO numerable, es decir no es finito y tampoco es numerable.}
\[
\mathcal C:=\{f:\mathbb N\to\{0,1\}\}.
\]
Entonces $\mathcal C$ es no numerable.

\section*{Foto 19}

Contin\'uamos con las definiciones.

\[
\mathcal C:=\{f:\mathbb N\to\{0,1\}\}.
\]
Sea $E\subset \mathcal C$ cualquier subconjunto numerable de $\mathcal C$.
P.D. que $E$ es un subconjunto propio.

\section*{Foto 20}

$\mathcal C$ es no numerable.

Espacios m\'etricos.

\textit{Definici\'on.}
Sea $X$ un conjunto no vac\'io. Una m\'etrica en $X$ es una funci\'on o distancia
\[
d:X\times X\to[0,\infty)
\]
con las siguientes propiedades:
\begin{enumerate}[label=\roman*)]
    \item $d(p,q)>0$ si $p\neq q$ y $d(p,p)=0$;
    \item $d(p,q)=d(q,p)$;
    \item $d(p,r)\le d(p,q)+d(q,r)$.
\end{enumerate}

\section*{Foto 21}

\textit{Ejemplos.}

\begin{enumerate}[label=\arabic*.]
    \item En $\mathbb R$, $d(a,b)=|a-b|$.
    \item En $\mathbb R^d$, $d(a,b)=\sqrt{\sum_{i=1}^d(a_i-b_i)^2}$.
    \item M\'etrica discreta:
    \[
    d(a,b)=
    \begin{cases}
    1,& a\neq b,\\
    0,& a=b.
    \end{cases}
    \]
\end{enumerate}

\section*{Foto 22}

Lema: la m\'etrica discreta es, en efecto, una m\'etrica en $X$.

\section*{Foto 23}

\textit{Ejemplo.} Sea $a,b\in\mathbb R$, $a<b$.

\[
C([a,b])=\{f:[a,b]\to\mathbb R\mid f \text{ es continua}\}.
\]

Sea
\[
d:C([a,b])\times C([a,b])\to[0,\infty),
\qquad
d(f,g)=\sup\{|f(x)-g(x)|\mid x\in[a,b]\}.
\]

\textit{Afirmaci\'on.} $d$ es m\'etrica en $C([a,b])$.

\section*{Foto 24}

\[
|f(x)-g(x)|\le |f(x)-h(x)|+|h(x)-g(x)|
\]
para todo $x\in[a,b]$.
Entonces
\[
d(f,g)\le d(f,h)+d(h,g).
\]

\section*{Foto 25}

\[
(C([a,b]),d)\text{ es un espacio m\'etrico.}
\]

\end{document}
